\chapter{Programas de ordenador (Categoría 6)}

\begin{introcategoria}
El código fuente va como \textbf{fichero adjunto (ZIP o tar.gz)} en este mismo depósito
y además depositado en Software Heritage. Aquí se describen los programas y en qué temas
se usan. \emph{Ojo:} una guía o un PDF con enlaces a repositorios \textbf{no} cuenta como
software (motivo de exclusión en convocatorias anteriores).
\end{introcategoria}

\section{Descripción de los programas}
\begin{tabularx}{\textwidth}{@{}l X l@{}}
\toprule
\textbf{Programa} & \textbf{Qué contiene y cómo se usa} & \textbf{Temas}\\
\midrule
practica01 & Plantilla y solución de la práctica 1 \ldots & 1, 2\\
\bottomrule
\end{tabularx}

\bigskip
\noindent Repositorio público: \url{https://github.com/usuario/repositorio}\\
SWHID: \texttt{swh:1:dir:xxxxxxxxxxxxxxxxxxxxxxxxxxxxxxxxxxxxxxxx}\\
Fichero adjunto en el depósito: \texttt{codigo-fuente.zip}

% Opcional: incluir el README en PDF
\material{README del repositorio}{materiales/6-software/README.pdf}{Instrucciones de uso del software.}
